Super-horizon freezing preserves the primordial amplitude of adiabatic curvature perturbations as a constant of the motion; the subsequent horizons density/pull constitutes the causal mechanism that converts that frozen amplitude into acoustic peaks in the cosmic microwave background, hierarchical density growth that seeds large-scale structure, and a geometrically calibrated elevation of the local expansion rate. When this causal sequence is realized on a discrete lattice whose scale-factor proxy advances as \(y_n=n\cdot\varepsilon\) with \(\varepsilon=10^{-9}\), every stage becomes numerically exact, topologically closed, and observationally falsifiable.

Step 1. Quantum Origin and Exit from the Horizon  
During inflation, quantum fluctuations of the inflaton generate a nearly scale-invariant spectrum of comoving curvature perturbations \(\mathcal{R}_k\). While a mode remains deep inside the Hubble radius its amplitude oscillates; once \(k\ll aH\) the spatial-gradient term in the Mukhanov–Sasaki equation becomes negligible and \(\mathcal{R}_k\) freezes to a constant value for adiabatic modes. This freezing is the geometric expression of causal disconnection: the mode’s wavelength exceeds the distance over which information can propagate, so its amplitude is locked relative to the expanding background.

Step 2. The Timespeed Frame Delay  
On the discrete lattice the same freezing appears as a prolonged interval during which the complex state vector \(z_n=y_n+iw_n\) advances only in its real (scale-factor) component while the phase remains stationary. The timespeed frame delay is therefore the lattice counterpart of super-horizon conservation: local proper-time ticks continue, yet the curvature mode experiences no evolution until the horizon condition is again satisfied.

Step 3. Horizon Re-entry and the Density/Pull Mechanism  
As the Hubble radius grows, each frozen mode eventually satisfies \(k\sim aH\). At that instant the previously negligible gradient terms reassert themselves. The gravitational potential \(\Phi\) (related to \(\mathcal{R}_k\) by \(\Phi\approx-\mathcal{R}_k\) on super-horizon scales and \(\Phi\approx\frac35\mathcal{R}_k\) in the matter era) begins to source fluid motions. This is the horizons density/pull: the horizon surface supplies the causal boundary, while the density contrast is gravitationally drawn toward overdense regions. The transfer function \(T(k,a)\) quantifies the conversion,
\[
\mathcal{P}(k,a)=\mathcal{P}_\mathcal{R}(k)\,T^2(k,a).
\]

Step 4. Conversion into Acoustic Peaks  
In the radiation-dominated plasma the re-entering modes drive coherent acoustic oscillations of the photon-baryon fluid. The frozen amplitude \(\mathcal{R}_k\) sets the initial conditions for these oscillations; the subsequent compression and rarefaction phases imprint the characteristic series of peaks and troughs observed in the CMB temperature and polarization spectra. Because the lattice preserves the exact value of \(\mathcal{R}_k\) until the discrete horizon-entry step, the acoustic-peak locations and amplitudes remain identical to those predicted by the continuous theory, up to the resolution of the grid \(\varepsilon\).

Step 5. Conversion into Density Growth  
After matter-radiation equality the same modes, now inside the horizon, source the growth of cold-dark-matter density contrasts \(\delta_m\propto a\). The horizons density/pull therefore translates the primordial amplitude into the hierarchical clustering that forms galaxies and the cosmic web. On the lattice this growth is realized by updating the real-space density field at each post-entry step while the topological winding (bosonic full-cycle or fermionic half-turn) remains closed, guaranteeing that no spurious phase noise contaminates the power spectrum.

Step 6. Calibration of the Local Expansion Rate  
At late times (\(z\lesssim0.15\)) the same causal sequence generates a residual geometric torque. The torque is confined by explicit suppression functions so that high-redshift observables (CMB, BAO) remain essentially untouched, while the local Hubble parameter is elevated from a baseline of \(70\) to \(73.17\,\mathrm{km\,s^{-1}\,Mpc^{-1}}\). Thus the horizons density/pull, acting on the previously frozen amplitude, simultaneously produces structure and resolves the Hubble tension without additional free parameters.

Step 7. Early-Universe Consistency on the Lattice  
The discrete evolution begins at the earliest epochs. The accompanying table records the neutron fraction \(X_n\) and helium-4 mass fraction \(Y_{\mathrm{He4}}\) at successive lattice points while the temperature remains \(T=1\,\mathrm{MeV}\):

| \(n\) | \(y_n\)     | \(T\) (MeV) | \(X_n\)  | \(Y_{\mathrm{He4}}\) |
|-------|-------------|-------------|---------|----------------------|
| 20    | \(2.0\times10^{-8}\) | 1.0000     | 0.4987 | 0.0                  |
| 100   | \(1.0\times10^{-7}\) | 1.0000     | 0.4974 | \(1.10\times10^{-31}\) |
| \ldots | \ldots     | 1.0000     | \ldots | \ldots               |
| 580   | \(5.8\times10^{-7}\) | 1.0000     | 0.4896 | \(3.10\times10^{-30}\) |

These values lie in the pre-BBN weak-interaction regime. The slow decline of \(X_n\) and the still-negligible \(Y_{\mathrm{He4}}\) demonstrate that the lattice correctly maintains thermal equilibrium and causal connectivity at temperatures where horizon scales are microscopic. Consequently the same discrete pathway that later freezes and re-enters the curvature modes is already consistent with the microphysical initial conditions required for successful Big-Bang nucleosynthesis.

Step 8. Closure and Falsifiability  
Because every stage—freezing, re-entry, acoustic driving, density growth, and late-time torque—is executed on a single, topologically closed lattice whose terminal monodromies return the system to the real axis with vanishing phase debt, the entire causal chain is both computationally reproducible and observationally accountable. Acoustic-peak positions, the matter power spectrum, the local \(H_0\), and the primordial helium abundance can all be extracted from the same run and compared directly with data. Simultaneous agreement across all observables constitutes a non-trivial confirmation of the thesis that super-horizon freezing followed by horizons density/pull is the unique causal mechanism linking the quantum origin of structure to the calibrated expansion rate of the present universe.
